\documentclass[11pt]{article}
\usepackage[margin=1in]{geometry}

% Core packages
\usepackage{amsmath,amssymb}
% \usepackage{tikz-cd}
\usepackage{multicol}
\usepackage{array}
\usepackage{booktabs}
\usepackage{longtable}
\usepackage{enumitem}
\usepackage{graphicx} % To include figures (Gantt, Pie chart, etc.)
\usepackage[hidelinks]{hyperref}

\renewcommand{\tablename}{Tabla}
\renewcommand{\contentsname}{Índice}
\renewcommand{\refname}{Bibliografía}
\renewcommand{\figurename}{Figura}

% Paragraphs
\setlength{\parindent}{0pt}
\setlength{\parskip}{1\baselineskip}

\title{Informe del Proyecto Programado\\Biblioteca UCR Recinto de Grecia}
\author{Marcos Ferreto Estrada \and Paulo Anchía Correas}
\date{\today}

\begin{document}

\begin{titlepage}
    \centering
    {\Large Universidad de Costa Rica\par}
    {\large Sede de Occidente - Recinto de Grecia\par}
    {\large Curso: Estructuras de Datos\par}

    \vspace{2.5cm}
    {\LARGE\bfseries Proyecto Programado\par}
    {\Large Sistema de gestión bibliotecaria\par}

    \vspace{2.5cm}
    \begin{tabular}{rl}
        Integrantes: & Marcos Ferreto Estrada --- C5F092 \\
                     & Paulo Anchía Correás --- C5C482 \\
        Profesora:   & Iyubanit Rodríguez Ramírez \\
        Fecha:       & 26 de junio del 2026 \\
    \end{tabular}

    \vfill
\end{titlepage}

\tableofcontents
\newpage

\section{Introducción}

Este documento describe casi todo el desarrollo del proyecto programado para administrar los datos de la biblioteca de la Universidad de Costa Rica, Recinto de Grecia. El sistema permite cargar, almacenar, consultar, eliminar y prestar libros mediante el uso de estructuras de datos que fueron estudiadas en el curso. 

\section{Descripción del problema}

El problema consiste en diseñar e implementar una aplicación que gestione información de libros, estudiantes y préstamos de una biblioteca utilizando estructuras de datos. El sistema permite buscar, insertar, eliminar y registrar préstamos, haciendo una solución integral a las necesidades de la biblioteca.

Todos los datos se quedan en archivos de texto plano delimitados (con extensión \texttt{.xml}) y, al ejecutar el programa se cargan en las estructuras de datos correspondientes en memoria principal. El sistema maneja libros identificados por un código único, estudiantes por su carné y préstamos con un código autogenerado. Además de que el programa incluye validaciones de negocio estrictas, impidiendo operaciones inválidas como prestar un libro que ya se encuentra prestado, o eliminar a un estudiante que mantenga devoluciones pendientes. Todo esto con la ayuda de una Interfaz Gráfica de Usuario (GUI) hecha con \texttt{tkinter} para facilitar la interacción visual y el uso del programa final.

\section{Análisis del problema}

\subsection{Estructuras de datos utilizadas}

\subsubsection{Árbol AVL para libros}

Los libros se almacenan en un árbol AVL ordenado por su código numérico único. Este método resulta mejor debido a que mantiene un balance automáticamente tras cada inserción o eliminación, asegurando un tiempo de complejidad de $O(\log n)$ en búsquedas, lo cual es importante cuando el catálogo de libros crece.

\textbf{Datos almacenados por libro:}
\begin{itemize}[noitemsep]
    \item Código único.
    \item Autor.
    \item Título.
    \item Año de publicación.
    \item Editorial.
    \item Área o temática.
\end{itemize}

\subsubsection{Tabla hash para estudiantes}

Los estudiantes se gestionan a través de una tabla hash. La función lo calcula a partir del número de carné de cada estudiante utilizando el módulo sobre el tamaño de la tabla o método de división.

Las colisiones se resuelven utilizando el encadenamiento a través de listas simples enlazadas. Esta estructura reduce el tiempo de acceso para ubicar el perfil de un estudiante.

\textbf{Datos almacenados por estudiante:}
\begin{itemize}[noitemsep]
    \item Carné único.
    \item Nombre.
    \item Carrera que cursa.
    \item Teléfono.
    \item Correo electrónico.
    \item Dirección física.
\end{itemize}

\subsubsection{Árbol rojinegro para préstamos}

Los registros de préstamos de libros se indexan en un Árbol Rojinegro. Al tratarse de operaciones transaccionales donde las inserciones y las consultas son frecuentes, esta estructura maneja muy bien el balance sin generar rotaciones excesivas (a diferencia del AVL). Mantiene los tiempos de operación acotados en $O(\log n)$ sin degradarse ante la inserción de fechas u órdenes consecutivos.

\textbf{Datos almacenados por préstamo:}
\begin{itemize}[noitemsep]
    \item Código único de préstamo.
    \item Código del libro prestado.
    \item Carné del estudiante solicitante.
    \item Fecha de préstamo.
    \item Fecha de devolución esperada.
\end{itemize}

\subsection{Carga y persistencia de archivos}

El programa cuenta con un módulo \texttt{XMLManager} especializado en la lectura de los archivos de texto \texttt{libros.xml}, \texttt{estudiantes.xml} y \texttt{prestamos.xml}, que pese a su nombre, almacenan atributos delimitados por el carácter de porcentaje (\%). Al iniciar la aplicación, se abren estos archivos y cada línea se separa para instanciar directamente los objetos \texttt{Libro}, \texttt{Estudiante} y \texttt{Prestamo}. Posteriormente, al momento de registrar un nuevo préstamo o al ejecutar una eliminación de un estudiante o libro válido, el programa recorre la estructura de datos correspondiente (haciendo uso del recorrido inorden o listando el hash) y sobrescribe los archivos, asegurando la persistencia de los cambios.

\subsection{Resolución de subproblemas}

El desarrollo del sistema se abordó de forma incremental, dividiendo el problema general en distintas etapas de desarrollo. Cada subproblema fue resuelto mediante una estructura de datos específica y validaciones de negocio acordes con los requisitos del sistema, siguiendo una implementación lógica y progresiva.

En la siguiente tabla se detalla cómo se abordó la implementación de cada requerimiento por etapas.

\begin{longtable}{p{0.15\textwidth} p{0.25\textwidth} p{0.50\textwidth}}
\toprule
\textbf{Etapa} & \textbf{Subproblema} & \textbf{Solución implementada} \\
\midrule

Etapa 1 
& Estructura base & Se configuró el proyecto inicial, creando los directorios principales y estableciendo las bases para los distintos módulos. \\
\midrule

Etapa 2 
& Modelos de datos & Se programaron las clases \texttt{Libro}, \texttt{Estudiante} y \texttt{Prestamo} con sus respectivos atributos para representar la información de manera estructurada en el sistema. \\
\midrule

Etapa 3 
& Persistencia XML & Se implementó un módulo para leer y escribir en archivos de texto, permitiendo cargar los datos al iniciar y guardar los cambios al finalizar una operación. \\
\midrule

Etapa 4
& Buscar libros por autor & Se realiza un recorrido inorden sobre el árbol AVL, evaluando si el autor del nodo actual coincide o contiene la subcadena ingresada. Cuando existe coincidencia, el libro se agrega a una lista de resultados. \\
\cmidrule{2-3}
& Buscar libro por título & Se aplica un procedimiento similar al de búsqueda por autor, recorriendo el AVL y comparando el título de cada nodo sin distinguir mayúsculas ni minúsculas. \\
\cmidrule{2-3}
& Buscar libro por código & Se aprovecha la propiedad de ordenamiento del AVL. La búsqueda compara recursivamente el código consultado con el código del nodo actual, reduciendo el tiempo de acceso a $O(\log n)$. \\
\cmidrule{2-3}
& Mostrar árboles en inorden & Se implementaron recorridos recursivos que visitan el hijo izquierdo, procesan la raíz y luego el hijo derecho, permitiendo mostrar los elementos en orden ascendente por código. \\
\midrule

Etapa 5
& Buscar estudiante por carné & Se calcula la función hash del carné para obtener el índice correspondiente en la tabla. Luego, se recorre la lista enlazada almacenada en esa posición hasta encontrar el registro exacto. \\
\cmidrule{2-3}
& Buscar estudiante por nombre & Se recorren todas las listas enlazadas de la tabla hash para encontrar coincidencias totales o parciales con el nombre ingresado. \\
\cmidrule{2-3}
& Buscar estudiantes por carrera & Se sigue un procedimiento similar al de búsqueda por nombre, recorriendo toda la tabla y acumulando los registros que coinciden con la carrera consultada. \\
\midrule

Etapa 6
& Estructura Rojinegra & Se implementó el árbol rojinegro con sus respectivas rotaciones e inserciones balanceadas para almacenar de manera eficiente los préstamos registrados. \\
\midrule

Etapa 7
& Sistema de préstamos & Primero se verifica la existencia del libro y del estudiante. Luego se comprueba que el libro no esté prestado. Si todo es correcto, se crea el préstamo con duración de 15 días, se inserta en el árbol rojinegro y se actualiza el archivo. \\
\midrule

Etapa 8
& Eliminar libro por código & Antes de eliminar, se valida en el árbol rojinegro que el libro no esté asociado a un préstamo activo. Si el libro está prestado, la eliminación se rechaza. En caso contrario, se elimina del AVL y se rebalancea el árbol. \\
\cmidrule{2-3}
& Eliminar estudiante por carné & Antes de eliminar, el sistema verifica que el estudiante no tenga préstamos activos en el árbol rojinegro. Si existe al menos uno, la eliminación se bloquea. Si no, se elimina el nodo correspondiente de la tabla hash. \\
\midrule

Etapa 9
& Interfaz gráfica & Se desarrolló con \texttt{tkinter}, organizando la aplicación en pestañas mediante \texttt{Notebook} para separar libros, estudiantes y préstamos. Además, se incorporaron formularios de entrada, botones de acción y un \texttt{Treeview} tabular para mostrar resultados. \\
\midrule

Etapa 10
& Documentación interna & Se agregaron \textit{docstrings} y comentarios descriptivos en el código, detallando los parámetros, retornos y el funcionamiento algorítmico de cada función. \\
\bottomrule
\end{longtable}

\subsection{Funciones o métodos principales}

A continuación, se listan los métodos de integración y gestión más importantes, correspondientes a los controladores del proyecto (\texttt{SistemaPrestamos}, \texttt{GestorEliminacion} y \texttt{XMLManager}).

\begin{longtable}{p{0.25\textwidth} p{0.24\textwidth} p{0.18\textwidth} p{0.23\textwidth}}
\toprule
\textbf{Función o método} & \textbf{Parámetros} & \textbf{Retorna} & \textbf{Descripción} \\
\midrule
\texttt{cargar\_todo} & Ninguno (lee de los atributos de clase) & Tupla de 3 listas de objetos & Lee las líneas de los archivos \texttt{.xml}, instancia cada clase y retorna los conjuntos en crudo para popular la GUI. \\
\texttt{dar\_prestamo} & \texttt{codigo\_libro}, \texttt{carnet} & \texttt{(bool, str)} & Valida y registra un préstamo si el libro está disponible. Retorna una tupla con un valor booleano de éxito y un mensaje de estado. \\
\texttt{eliminar\_libro} & \texttt{codigo} & \texttt{(bool, str)} & Realiza las verificaciones de préstamos activos y elimina el libro del AVL si todo es correcto, luego sincroniza en disco. \\
\texttt{eliminar\_estudiante}& \texttt{carnet} & \texttt{(bool, str)} & Ejecuta validaciones análogas pero dirigidas a la tabla hash del estudiante, también sincronizando al disco en caso exitoso. \\
\bottomrule
\end{longtable}

\section{Análisis de resultados}

La siguiente tabla resume el grado de completitud técnica de las diferentes etapas y requerimientos del sistema en el momento de la entrega del código final.

\begin{longtable}{p{0.38\textwidth} p{0.24\textwidth} p{0.28\textwidth}}
\toprule
\textbf{Parte del proyecto} & \textbf{Estado} & \textbf{Observaciones} \\
\midrule
Archivos XML & Terminado con éxito & Uso correcto de la persistencia directa mediante porcentaje (\%). \\
Árbol AVL & Terminado con éxito & Funcional con rebalanceo automático, inserción, búsqueda y eliminación. \\
Tabla hash & Terminado con éxito & Funcional con resolución de colisiones y extracción de datos. \\
Árbol rojinegro & Terminado con éxito & Funcional en balanceo por colores para administrar los registros. \\
Búsquedas & Terminado con éxito & Algoritmos de búsqueda acoplados correctamente a la GUI. \\
Carga de archivos & Terminado con éxito & Funciona eficientemente para poblar los árboles al iniciar. \\
Eliminación de datos & Terminado con éxito & Validaciones integradas que impiden corromper los préstamos. \\
Préstamos & Terminado con éxito & Validación lógica estricta y control de fechas. \\
Interfaz gráfica de usuario & Terminado con éxito & Pantallas organizadas por pestañas, componentes limpios. \\
Documentación interna & Terminado con éxito & Funciones documentadas e informe elaborado. \\
\bottomrule
\end{longtable}

\section{Conclusiones y recomendaciones}

\begin{enumerate}
    \item \textbf{Conclusión:} El uso adecuado de las estructuras de datos mejoró sustancialmente la eficiencia y escalabilidad del sistema comparado a utilizar simples arreglos secuenciales.\\
    \textbf{Recomendación:} Es recomendable mantener un análisis de complejidad inicial antes de empezar futuros desarrollos para escoger el modelo de datos correcto desde la fase de planeación.

    \item \textbf{Conclusión:} El Árbol AVL demostró ser altamente útil y óptimo para mantener los códigos de libros ordenados y responder rápidamente a búsquedas exactas.\\
    \textbf{Recomendación:} En implementaciones donde las eliminaciones de libros sean abrumadoramente masivas, el re-balanceo constante podría penalizar el rendimiento, por lo que podría recomendarse estudiar alternativas de baja frecuencia en ese caso hipotético.

    \item \textbf{Conclusión:} La tabla hash permitió el acceso directo y ágil al registro de los estudiantes mediante su número de carné, optimizando significativamente la operación.\\
    \textbf{Recomendación:} Para bases de datos universitarias reales masivas, se recomienda calibrar y testear el tamaño de la tabla (actualmente 100) frente al uso esperado para no castigar el factor de carga ni generar listas de colisión excesivamente largas.

    \item \textbf{Conclusión:} El manejo del Árbol Rojinegro brindó la estabilidad necesaria para las operaciones de préstamos, las cuales usualmente implican un crecimiento secuencial del identificador.\\
    \textbf{Recomendación:} Aprovechar el ordenamiento del árbol rojinegro para en una futura iteración añadir reportes automáticos de libros próximos a vencer su fecha de devolución.

    \item \textbf{Conclusión:} La programación de las validaciones de dependencia previas (comprobar si hay un préstamo antes de permitir borrar un libro) previno proactivamente las inconsistencias en disco e interrupciones del programa.\\
    \textbf{Recomendación:} Mantener una lógica fuerte del lado del controlador de persistencia antes de alterar memoria, para que sirva de segunda capa de validación.

    \item \textbf{Conclusión:} La Interfaz Gráfica de Usuario (GUI) construida con los componentes nativos, le dio un aspecto intuitivo y organizado a una lógica subyacente que de otra forma sería compleja de probar por comandos.\\
    \textbf{Recomendación:} Se puede ampliar el sistema incluyendo filtros avanzados por fechas, reportes PDF o alertas visuales cuando un estudiante acumula penalizaciones y morosidades, lo cual complementaría excelentemente a la GUI construida.

    \item \textbf{Conclusión:} El manejo de notificaciones se implementó de forma muy robusta mediante una barra de estado en el pie de página (footer) de la interfaz. Esto asegura que se muestre un error explícito si faltó algún dato, o un mensaje de éxito confirmando exactamente qué registro se insertó, buscó o eliminó al interactuar con el sistema.\\
    \textbf{Recomendación:} Conservar y escalar este patrón de diseño en futuros módulos de la aplicación, ya que ofrece una excelente retroalimentación en tiempo real al usuario sin interrumpir su flujo de trabajo con ventanas invasivas.
\end{enumerate}

\section{Bibliografía}

\begin{itemize}
    \item Python Software Foundation. (2026). \textit{tkinter — Interfaz gráfica de usuario}. Documentación oficial de Python. Recuperado de \url{https://docs.python.org/es/3/library/tkinter.html}
    \item Cormen, T. H., Leiserson, C. E., Rivest, R. L., \& Stein, C. (2009). \textit{Introduction to Algorithms} (3rd ed.). MIT Press.
    \item Material didáctico y diapositivas de la cátedra del curso "Estructuras de Datos", Universidad de Costa Rica, Recinto de Grecia.
\end{itemize}

\end{document}
